\uuid{Fj8x}
\titre{Module/argument}
\auteur{bodin ; hanani ; mzari}
\organisation{exo7}
\licence{CC BY-NC-SA 4.0}
\datecreate{2020-01-20}
\niveau{L1}
\module{Algèbre}
\chapitre{Nombres complexes}
\sousChapitre{Forme cartésienne, forme polaire}
\difficulte{2}
\theme{Écritures algébrique et géométrique, module, argument}
\exo7qcmid{lille-1.nombres-complexes.009}

\contenu{
\begin{question}
\qclassification{104.01}
\qtags{motcle=module/argument}
Soit $z_1$ et $z_2$ deux nombres complexes. Alors, $|z_1+z_2|^2 + |z_1-z_2|^2$ est égal à :
\begin{answers}
    \bad{$|z_1|^2+|z_2|^2$}

    \bad{$|z_1|^2-|z_2|^2$}

    \good{$ 2|z_1|^2 +2|z_2|^2$}

    \bad{$ 2|z_1|^2 -2|z_2|^2$}

\end{answers}
\begin{explanations}
Utiliser :  $|z|^2= z\overline{z}$.
\end{explanations}
\end{question}
}
